% Verbatim from minor_report/chapters/chapter_3_methodology.tex lines 164-298 (retired report),
% headings demoted by 0 level(s). Do not edit here; edit the source of the numbers.

The simulated Ackermann robot carries two ideal pinhole cameras and an IMU on a
forward mast, and is driven through a front-steering, rear-driven four-wheel
chassis. \Cref{fig:ackermann-robot} gives its dimensioned plan and elevation, and
\cref{tab:sim-sensor-rig,tab:platform-geometry} list the sensor and platform
constants. Every value in the figure and both tables is read from
\repo{aws-robomaker-small-warehouse-world/models/ackermann_robot/model.sdf}, so
the drawing is to scale rather than illustrative.

\begin{figure}[htbp]
  \centering
  \resizebox{\textwidth}{!}{% Dimensioned plan and elevation of the simulated Ackermann robot.
%
% Every coordinate below is read from
% aws-robomaker-small-warehouse-world/models/ackermann_robot/model.sdf and is
% expressed in metres relative to base_footprint, so the drawing is to scale and
% can be re-derived from the model rather than trusted. A Gazebo screenshot was
% deliberately not used: the quantities that matter to the estimators are the
% wheel base, track, wheel radius and sensor offsets, and a perspective render
% shows none of them.
%
% NOTE ON THE TWO VIEWS. They are placed side by side with shift={(1.45,0)},
% i.e. in DIAGRAM units, not with yshift/xshift in centimetres. Inside a scope
% that already carries scale=9 a shift given in cm is multiplied by that scale,
% which silently pushes the second view about 56 cm away and off the page.
\begin{tikzpicture}[
  scale=9,
  chassis/.style={draw, fill=gray!12, thick},
  wheel/.style={draw, fill=gray!45, thick},
  steered/.style={draw, dashed, gray!60!black, line width=0.4pt},
  mast/.style={draw, fill=gray!25},
  sensor/.style={draw, fill=external, thick},
  imu/.style={draw, fill=implemented, thick},
  dim/.style={{Latex[length=1.2mm]}-{Latex[length=1.2mm]}, gray!60!black,
    line width=0.3pt},
  dimlab/.style={font=\tiny, inner sep=1pt, fill=white},
  leader/.style={gray!60!black, line width=0.25pt},
  lbl/.style={font=\tiny, inner sep=1pt},
  tick/.style={gray!60!black, line width=0.25pt}
]

% ======================= PLAN VIEW (x forward, y left) =======================
\begin{scope}
  \node[font=\small\bfseries] at (0.15,0.44) {Plan view};

  % chassis deck: 0.6 x 0.463 centred on base_footprint
  \draw[chassis] (-0.3,-0.2315) rectangle (0.3,0.2315);

  % steered front wheels shown at full lock, +/-0.6109 rad = +/-35 deg,
  % rotated about the kingpin rather than drawn as a bare angle sweep
  \foreach \y in {0.21,-0.21}{
    \begin{scope}[shift={(0.21,\y)}, rotate=35]
      \draw[steered] (-0.0585,-0.0325) rectangle (0.0585,0.0325);
    \end{scope}
  }

  % wheels: radius 0.0585, half-width 0.0325, at x = +/-0.21, y = +/-0.21
  \foreach \x in {0.21,-0.21}{
    \foreach \y in {0.21,-0.21}{
      \draw[wheel] (\x-0.0585,\y-0.0325) rectangle (\x+0.0585,\y+0.0325);
    }
  }
  \node[lbl, anchor=south west] at (0.135,0.275)
    {$\delta_{\max}=\pm\SI{0.6109}{\radian}$};

  % camera mast and the two cameras
  \draw[mast] (0.3078,-0.01) rectangle (0.5395,0.01);
  \draw[sensor] (0.486,0.0836) rectangle (0.506,0.1036);
  \draw[sensor] (0.486,-0.01) rectangle (0.506,0.01);

  % IMU
  \draw[imu] (0.313,-0.025) rectangle (0.363,0.025);
  \node[lbl, anchor=north] at (0.338,-0.032) {\texttt{imu}};

  % base_footprint origin marker
  \draw[thick] (-0.025,0) -- (0.025,0) (0,-0.025) -- (0,0.025);
  \node[lbl, anchor=north] at (0,-0.035) {\texttt{base\_footprint}};

  % camera labels, lifted clear on leaders so they cannot collide with the
  % baseline dimension between them
  \draw[leader] (0.506,0.0936) -- (0.60,0.155);
  \node[lbl, anchor=west] at (0.605,0.155) {\texttt{cam0} (left)};
  \draw[leader] (0.506,0) -- (0.60,-0.075);
  \node[lbl, anchor=west] at (0.605,-0.075) {\texttt{cam1} (right)};

  % --- dimensions ---
  \draw[tick] (0.21,-0.2425) -- (0.21,-0.315);
  \draw[tick] (-0.21,-0.2425) -- (-0.21,-0.315);
  \draw[dim] (-0.21,-0.30) -- (0.21,-0.30);
  \node[dimlab] at (0,-0.30) {wheel base \SI{0.42}{\metre}};

  \draw[tick] (-0.2685,0.21) -- (-0.375,0.21);
  \draw[tick] (-0.2685,-0.21) -- (-0.375,-0.21);
  \draw[dim] (-0.36,-0.21) -- (-0.36,0.21);
  \node[dimlab, rotate=90] at (-0.36,0) {track \SI{0.42}{\metre}};

  \draw[tick] (0.506,0.0936) -- (0.565,0.0936);
  \draw[tick] (0.506,0) -- (0.565,0);
  \draw[dim] (0.552,0) -- (0.552,0.0936);
  \node[dimlab, rotate=90] at (0.552,0.047) {\SI{0.0936}{\metre}};
\end{scope}

% ======================= ELEVATION (x forward, z up) =======================
\begin{scope}[shift={(1.45,-0.16)}]
  \node[font=\small\bfseries] at (0.15,0.60) {Side elevation};

  % ground line
  \draw[gray!55, line width=0.5pt] (-0.38,0) -- (0.60,0);

  % wheels seen edge on: centre z = 0.072, radius 0.0585
  \foreach \x in {0.21,-0.21}{
    \draw[wheel] (\x,0.072) circle (0.0585);
    \draw[gray!60!black, line width=0.25pt] (\x,0.072) circle (0.007);
  }

  % chassis deck: z from 0.072 to 0.122
  \draw[chassis] (-0.3,0.072) rectangle (0.3,0.122);

  % mast: vertical post, camera bar, short riser
  \draw[mast] (0.399,0.122) rectangle (0.419,0.3715);
  \draw[mast] (0.3078,0.3715) rectangle (0.5395,0.3915);
  \draw[mast] (0.486,0.3915) rectangle (0.506,0.4295);

  % cameras at z = 0.4395 (cam0 hides cam1 in this projection)
  \draw[sensor] (0.482,0.4295) rectangle (0.510,0.4495);
  \draw[leader] (0.510,0.4395) -- (0.55,0.50);
  \node[lbl, anchor=west] at (0.555,0.50) {\texttt{cam0}, \texttt{cam1}};

  % IMU at z = 0.192
  \draw[imu] (0.313,0.182) rectangle (0.363,0.202);
  \draw[leader] (0.363,0.192) -- (0.44,0.25);
  \node[lbl, anchor=west] at (0.445,0.25) {\texttt{imu}, $z=\SI{0.192}{\metre}$};

  % --- dimensions ---
  \draw[dim] (0.21,0.072) -- (0.21,0.1305);
  \node[dimlab, anchor=west] at (0.222,0.101) {$r=\SI{0.0585}{\metre}$};

  \draw[tick] (0.510,0.4395) -- (0.585,0.4395);
  \draw[dim] (0.572,0) -- (0.572,0.4395);
  \node[dimlab, rotate=90] at (0.572,0.22) {\SI{0.4395}{\metre}};

  \draw[tick] (-0.2685,0.072) -- (-0.375,0.072);
  \draw[dim] (-0.36,0) -- (-0.36,0.072);
  \node[dimlab, rotate=90] at (-0.36,0.036) {axle \SI{0.072}{\metre}};
\end{scope}

\end{tikzpicture}
}
  \caption[Simulated Ackermann platform, drawn to scale from the model definition.]{Simulated Ackermann platform, drawn to scale from the model
  definition. Distances are in metres relative to \texttt{base\_footprint}. The
  stereo pair is mounted asymmetrically: \texttt{cam1} sits on the vehicle
  centreline and \texttt{cam0} is offset \SI{0.0936}{\metre} to its left, so the
  optical centre of the pair is not the chassis centreline. The wheel base and
  track are equal at \SI{0.42}{\metre}, which is what makes the rear-wheel
  differential a poor yaw-rate estimator (\cref{sec:proprio-method}).}
  \label{fig:ackermann-robot}
\end{figure}

\begin{table}[htbp]
\centering
\caption[Simulated sensor configuration.]{Simulated sensor configuration. Intrinsics are analytic rather than
calibrated: the simulator renders an exact pinhole camera, so $f_x=f_y$ follow
from the image width and the horizontal field of view and the distortion
coefficients are identically zero. The stereo baseline was independently checked
from positive disparity in a recorded bag rather than taken from the model alone.}
\label{tab:sim-sensor-rig}
\small
\resizebox{\textwidth}{!}{%
\begin{tabular}{llll}
\toprule
Quantity & \texttt{cam0} (left) & \texttt{cam1} (right) & \texttt{imu} \\
\midrule
Resolution & $1280\times720$ & $1280\times720$ & -- \\
Rate & \SI{30}{\hertz} & \SI{30}{\hertz} & \SI{200}{\hertz} \\
Horizontal FOV & $\pi/2$ rad (\SI{90}{\degree}) & $\pi/2$ rad (\SI{90}{\degree}) & -- \\
$f_x=f_y$ & 640 & 640 & -- \\
$(c_x,c_y)$ & $(640,360)$ & $(640,360)$ & -- \\
Distortion & zero & zero & -- \\
Position from \texttt{base\_footprint} (m) & $(0.496,0.0936,0.4395)$ &
  $(0.496,0,0.4395)$ & $(0.338,0,0.192)$ \\
Position from IMU body (m) & $(0.158,0.0936,0.2475)$ & $(0.158,0,0.2475)$ &
  $(0,0,0)$ \\
Stereo baseline & \multicolumn{2}{l}{\SI{0.0936}{\metre}} & -- \\
\bottomrule
\end{tabular}}
\end{table}

The fixed camera rotation converts the ROS body convention (forward, left, up) to
the optical convention (right, down, forward). The VINS simulation configuration
fixes the extrinsics (\repo{estimate_extrinsic: 0}) and disables temporal-offset
estimation (\repo{estimate_td: 0}); the assumed inertial noise model is given in
\cref{tab:imu-noise-assumed}, and the same values are mapped to the ORB-SLAM3
configuration.

\begin{table}[htbp]
\centering
\caption[Inertial noise parameters supplied to the visual-inertial estimators, beside the values measured from the recorded IMU stream in \cref{sec:proprio-method}.]{Inertial noise parameters supplied to the visual-inertial estimators,
beside the values measured from the recorded IMU stream in
\cref{sec:proprio-method}. The configured values are inherited estimator defaults,
not measurements of this simulated IMU; the discrepancy is carried forward as a
queried item rather than silently reconciled.}
\label{tab:imu-noise-assumed}
\small
\begin{tabular}{llll}
\toprule
Parameter & Configured (VINS, ORB) & Measured from bag & Units \\
\midrule
\texttt{acc\_n} & 0.02 & \flag{0.0139--0.0294} &
  m\,s$^{-2}/\sqrt{\mathrm{Hz}}$ \\
\texttt{gyr\_n} & 0.002 & \flag{0.000300--0.000364} &
  rad\,s$^{-1}/\sqrt{\mathrm{Hz}}$ \\
\texttt{acc\_w} & 0.001 & not measurable & m\,s$^{-3}/\sqrt{\mathrm{Hz}}$ \\
\texttt{gyr\_w} & 0.0001 & not measurable &
  rad\,s$^{-2}/\sqrt{\mathrm{Hz}}$ \\
\bottomrule
\end{tabular}
\end{table}

\begin{queried}[Queried]
The configured gyroscope noise density is \num{2e-3} whereas the stationary
windows of the five recorded streams give \numrange{3.00e-4}{3.64e-4}, so the
configured value is about six times too large. The configured accelerometer
density of 0.02 sits inside the measured spread but that spread is itself wide,
\numrange{0.0139}{0.0294} across five bags of the same simulated sensor. Both
estimators are therefore running on an inertial noise model that does not
describe this IMU closely. The bias random walks cannot be checked at all: the
longest stationary interval in any \repo{dataset_allsensor} bag is
\SI{1.54}{\second} and an Allan-variance estimate needs minutes. No conclusion in
this report depends on these four numbers, but they should be re-derived from a
dedicated stationary recording before any tuning claim is made.
\end{queried}

\begin{queried}[Queried]
The accelerometer characterisation is not stable across recordings of the same
simulated sensor. \repo{allsensor_002}, whose stationary window is
\SI{1.54}{\second} and therefore the best conditioned of the three, returns an $x$
bias of \SI{0.0683}{\metre\per\second\squared}, against 0.0191 for
\repo{allsensor_000} and \SI{0.0016}{\metre\per\second\squared} for
\repo{allsensor_004}, whose window is the shortest at \SI{0.52}{\second}. The
scale ranges from 1.0199 to 1.0669 over the five. Because the window is short in
every case, part of this spread is estimator variance rather than a property of
the sensor, and the accelerometer terms should not be quoted as calibrated
constants. The gyroscope scale agrees to four decimal places across all five bags;
its bias does not, varying by a factor of twenty, which is consistent with the
same window-length problem.
\end{queried}

\begin{table}[htbp]
\centering
\caption[Platform geometry and drivetrain limits used by the wheel-odometry and EKF baselines.]{Platform geometry and drivetrain limits used by the wheel-odometry and
EKF baselines. \texttt{AckermannSteering} exposes no noise parameter, so the
simulated encoders and odometry are noise free at source.}
\label{tab:platform-geometry}
\small
\begin{tabular}{llll}
\toprule
Quantity & Symbol & Value & Source field \\
\midrule
Wheel base (front to rear axle) & $L$ & \SI{0.42}{\metre} & \texttt{wheel\_base} \\
Track (wheel separation) & $B$ & \SI{0.42}{\metre} & \texttt{wheel\_separation} \\
Kingpin width & -- & \SI{0.42}{\metre} & \texttt{kingpin\_width} \\
Wheel radius & $r$ & \SI{0.0585}{\metre} & \texttt{wheel\_radius} \\
Wheel width & -- & \SI{0.065}{\metre} & wheel \texttt{<length>} \\
Steering limit & $\delta_{\max}$ & $\pm\SI{0.6109}{\radian}$ & \texttt{steering\_limit} \\
Chassis mass & -- & \SI{10.1}{\kilogram} & \texttt{base\_footprint} \texttt{<mass>} \\
Maximum speed & -- & \SI{10}{\metre\per\second} & \texttt{max\_velocity} \\
Maximum acceleration & -- & \SI{3}{\metre\per\second\squared} & \texttt{max\_acceleration} \\
Odometry publish rate & -- & \SI{50}{\hertz} & \texttt{odom\_publish\_frequency} \\
Joint-state publish rate & -- & \SI{1}{\kilo\hertz} & measured from bag \\
\bottomrule
\end{tabular}
\end{table}
